\chapter{Implementation and Measurement}

% ─────────────────────────────────────────────
\section{Experimental Setup}

\subsection{Hardware Setup}
\label{sec:implementation_hardware}

A critical decision in the experimental design was the selection of the computing platform on which all experiments would be conducted. Two options were considered: the university's shared computing infrastructure or a local machine. The latter was ultimately chosen, as it provides several key advantages for this study. First, it ensures complete reproducibility, since all experiments are executed on identical hardware configuration, eliminating interfering variables that could arise from shared resource conflicts or system variability. Second, it provides consistent availability and direct control over system configuration, which is particularly important given the extended runtime of the large-scale experiments (particularly on the 11M-instance HIGGS dataset). Third, local execution facilitates fine-grained performance monitoring and measurement of energy consumption, a primary objective of this study.

The hardware platform selected for all experiments is a consumer-grade desktop system with the following specifications: an NVIDIA RTX 4070 graphics processing unit (12 GB GDDR6 memory, 5,888 CUDA cores), an AMD Ryzen 7 5700X processor (8 cores / 16 threads, base clock 3.6 GHz), and 32 GB of DDR4 RAM. The system runs Windows 11 with all code implemented in Python 3.12.0.

\begin{table}[htbp]
\centering
\caption{Hardware specification of the experimental system}
\label{tab:hardware}
\begin{tabular}{ll}
\toprule
\textbf{Component} & \textbf{Specification} \\
\midrule
CPU & AMD Ryzen 7 5700X (8 cores / 16 threads, 3.6 GHz base) \\
GPU & NVIDIA RTX 4070 (12 GB GDDR6, 5{,}888 CUDA cores) \\
RAM & 32 GB DDR4 \\
Operating System & Windows 11 \\
\bottomrule
\end{tabular}
\end{table}

The choice of hardware reflects a deliberate trade-off between experimental scope and practical feasibility. The RTX 4070 represents a mid-to-high-end consumer GPU with sufficient compute capability to execute GPU-accelerated training for XGBoost and MLP models in reasonable time, while remaining accessible to many practitioners. The Ryzen 7 5700X provides adequate CPU performance for baseline models (Logistic Regression, Random Forest) and CPU-based XGBoost training. Together, this configuration enables direct comparison of CPU- and GPU-accelerated training on the same platform, which is central to answering one of the research questions (\textit{How does energy consumption differ between CPU-based and GPU-based training for equivalent tasks?}).

It is important to note that energy consumption and CO2 emissions measured in this study are specific to this hardware configuration and the local electricity grid, and may not directly generalize to other systems or geographic regions. However, the relative performance comparisons (e.g., the energy efficiency ranking of different algorithms) are more likely to generalize, as they reflect algorithmic properties rather than hardware specifics alone.

\subsection{Software Environment}
\label{implementation_sw_env}

All experiments were implemented in Python 3.12.0. The main libraries used are described below.

Scikit-learn \cite{pedregosaScikitlearnMachineLearning2011} served as the foundation of the evaluation pipeline. It provided the \texttt{LogisticRegression} and \texttt{RandomForestClassifier} implementations, as well as \texttt{KFold} and \texttt{cross\_validate} for cross-validation, \texttt{StandardScaler} for feature normalization, \texttt{make\_pipeline} for chaining preprocessing and model steps, and \texttt{f1\_score} and \texttt{make\_scorer} for evaluation metrics.

The XGBoost library \cite{chenXGBoostScalableTree2016} was used for the gradient boosting model, as it provides a highly optimized implementation that supports both CPU and GPU training.

For the MLP, PyTorch \cite{paszkePyTorchImperativeStyle2019} was used as the deep learning backend. To integrate the PyTorch model into the scikit-learn evaluation pipeline, the Skorch library \cite{ghoshSkorchBridgePyTorch2024} was used, which wraps PyTorch models in a scikit-learn-compatible interface.

Hyperparameter tuning for all models except Logistic Regression was performed using Optuna, as described in \Cref{design_hyperparameter}.

Carbon emissions and energy consumption were tracked using CodeCarbon \cite{courtyMlco2CodecarbonV2412024}, which is described in detail in \Cref{implementation_emissions}.

To ensure reproducibility, a fixed random seed of 42 was used for all data splits, cross-validation folds, and model initialization. A \texttt{requirements.txt} file is provided so that the exact library versions used in this study can be reproduced.

\newpage
\section{Dataset Integration}
\label{implementation_datasets}

The three datasets used in this study — \textit{Wine}, \textit{Credit}, and \textit{HIGGS} — were introduced in \Cref{experimental_dataset}. This section describes how they are integrated into the software architecture, covering the central configuration file, the data loading pipeline, and dataset-specific constraints.

\subsection{Configuration Architecture}
\label{implementation_config}

To ensure a consistent foundation across all model scripts, a central configuration 
file (\texttt{config.py}) was created. It serves as a single source of truth for 
all dataset-related settings, so that each script accesses the same parameters 
without redundant definitions.

For each dataset, the configuration stores the file path, the name of the target 
column, and an optional row limit (\texttt{nrows}). The row limit was particularly 
relevant for the HIGGS dataset, which contains 11 million instances, allowing 
experiments to be run on smaller subsets where needed. Beyond these common fields, 
the configuration also captures dataset-specific properties. For example, the Credit 
dataset includes a header row that must be skipped during loading, and the Wine 
dataset does not contain column names in the source file, which are therefore defined 
in the configuration directly. Additionally, a label offset is applied to the Wine 
target column, as its class labels are encoded as 1, 2, and 3 rather than starting 
at 0, which is required by the classifiers used in this study.

The global constants \texttt{RANDOM\_STATE} and \texttt{CV\_FOLDS} are also defined 
here, ensuring that all scripts use identical values for random seeding and 
cross-validation, which is essential for the comparability of results.

\subsection{Data Loading Pipeline}
\label{implementation_data_loading}

To read the datasets from their respective files, a custom \texttt{load\_data} method was implemented. Depending on the dataset, the appropriate loading mechanism is selected. The \textit{Wine} and \textit{Credit} datasets are loaded using pandas' \texttt{read\_csv} function. For the \textit{Higgs} dataset, \texttt{read\_parquet} is utilized. This was a deliberate decision to save storage space by storing the large \textit{Higgs} dataset as a Parquet file rather than a standard CSV.

Following the initial loading phase, the \textit{Higgs} data is reduced to a representative subset using a sampling approach controlled by the \texttt{nrows} parameter. Once the data size is managed, any columns specified in the configuration file are dropped from the respective dataset. The data is then split into the feature matrix $X$ and the target vector $y$. If an offset is defined in the config, which is currently only the case for the \textit{Credit} dataset, it is applied as described in the previous section. Finally, the method returns $X$ and $y$ for further use in the subsequent scripts. This process is visualized in \Cref{fig:data_loading_pipeline}.

Due to the substantial scale of the \textit{Higgs} dataset, several constraints were implemented to ensure the stability of the execution environment. Training a \texttt{Random Forest} model on this specific data led to unmanageable processing times. For this reason, the \texttt{Random Forest} algorithm is explicitly bypassed for the \textit{Higgs} dataset and the program terminates using \texttt{sys.exit(0)} to prevent system crashes. This allows the process to stay within the boundaries of the available hardware resources.

\begin{figure}[htbp]
\centering
\begin{tikzpicture}[
    node distance=1.5cm and 2cm,
    process/.style={rectangle, draw, rounded corners, minimum width=3.5cm, minimum height=1cm, text centered, font=\sffamily},
    decision/.style={diamond, draw, text centered, aspect=2, inner sep=0pt, font=\sffamily},
    arrow/.style={thick, ->, >=stealth}
]
\node (dataset) [decision] {Dataset?};
\node (load_csv) [process, below left=1cm and 0.5cm of dataset] {\texttt{read\_csv()}};
\node (load_parq) [process, below right=1cm and 0.5cm of dataset] {\texttt{read\_parquet()}};
\node (sample) [process, below=1cm of load_parq] {Sample \texttt{nrows}};
\node (drop) [process, below=4.5cm of dataset] {Drop config columns};
\node (split) [process, below=1cm of drop] {Split $X$ and $y$};
\node (offset) [decision, below=1cm of split] {Offset defined?};
\node (apply) [process, right=1.5cm of offset] {Apply offset};
\node (return) [process, below=1cm of offset] {Return $X$ and $y$};

\draw [arrow] (dataset) -| node[anchor=east, yshift=0.2cm] {\textit{Wine} \& \textit{Credit}} (load_csv);
\draw [arrow] (dataset) -| node[anchor=west, yshift=0.2cm] {\textit{Higgs}} (load_parq);
\draw [arrow] (load_parq) -- (sample);
\draw [arrow] (load_csv) |- (drop);
\draw [arrow] (sample) |- (drop);
\draw [arrow] (drop) -- (split);
\draw [arrow] (split) -- (offset);
\draw [arrow] (offset) -- node[anchor=south] {Yes} (apply);
\draw [arrow] (offset) -- node[anchor=east] {No} (return);
\draw [arrow] (apply) |- (return);
\end{tikzpicture}
\caption{Visualization of the data loading pipeline and preprocessing steps within the \texttt{load\_data()} method.}
\label{fig:data_loading_pipeline}
\end{figure}

\section{Reproducibility and Evaluation Design}
\label{implementation_reproducibility}

To ensure fully reproducible results across all experiments, a global constant \texttt{RANDOM\_STATE} with a value of 42 is utilized in every model script. This constant is applied to all functions requiring stochastic consistency. It is specifically passed to the \texttt{KFold} splitter, the \texttt{train\_test\_split} method, and the initialization of the machine learning models. For the multilayer perceptron, the random seed is explicitly set using \path{torch.manual_seed(RANDOM_STATE)} to guarantee deterministic weight initialization and training behavior.

The evaluation process relies on a standardized cross validation setup implemented uniformly across all scripts. The data is partitioned using \texttt{KFold} configured with five splits, enabled shuffling, and the aforementioned random state. The execution is handled by the \texttt{cross\_validate()} function.

The performance metrics are evaluated using a custom scoring dictionary. This dictionary tracks both the standard accuracy and a weighted F1 score constructed via \texttt{make\_scorer(f1\_score, average="weighted")}. Once the evaluation is complete, the overall performance is aggregated by calculating the mean of the individual fold results, which is accessed through \texttt{cv\_results["test\_accuracy"].mean()} alongside the corresponding F1 score array. Having defined the measures of reproducibility and testing, the focus now shifts to the actual implementation of the predictive models.

\section{Model Implementations}
\label{implementation_models}

As outlined in \Cref{design_models}, the following classification algorithms were employed to address the research questions mentioned previously. This section therefore details the specific implementation of each model.

\subsection{Logistic Regression}
\label{implementation_lr}

The first baseline model is a standard logistic regression. To ensure the model processes all features on an equal scale, the algorithm is embedded in a pipeline starting with a \texttt{StandardScaler}. This initial scaling step is crucial for the stability and overall performance of the subsequent classification. 

The regression is configured to use the Stochastic Average Gradient (SAG) algorithm by setting the \texttt{solver} parameter to \texttt{sag}. This specific solver is highly optimized for large datasets. It significantly reduces the memory footprint and required training time, which is especially important when evaluating the extensive \textit{Higgs} dataset. To ensure proper convergence, the maximum number of iterations is capped at 1000 using the \texttt{max\_iter} parameter. 

This model setup deliberately excludes any hyperparameter optimization. Instead, it acts as a stable baseline to evaluate the raw computational time without the extra overhead of complex search processes.

\subsection{Random Forest}
\label{implementation_rf}

The second baseline algorithm is the Random Forest classifier. This method belongs to the family of ensemble learning techniques. Unlike logistic regression, this algorithm does not require prior feature scaling. To maximize computational efficiency, the execution is fully parallelized across all available CPU cores by setting the \texttt{n\_jobs} parameter to \texttt{-1}. This configuration allows the processor to run at maximum capacity, which directly influences the subsequent energy consumption tracking and overall runtime measurements.

An important methodological decision is the explicit exclusion of the \textit{Higgs} dataset for this specific classifier. As established in \Cref{implementation_data_loading}, training a Random Forest on such a massive volume of data leads to unmanageable runtimes. The primary research question focuses on the balance between energy consumption and predictive performance across varying dataset sizes and model complexities. Forcing this algorithm to process the \textit{Higgs} data would result in a massive energy footprint and an excessive runtime. This extreme resource consumption would severely skew the comparative baseline without providing proportional insights into the underlying scaling behavior. Therefore, this execution is bypassed to maintain a sensible and fair evaluation framework.

To ensure maximum predictive capabilities for the remaining data, the model configurations were optimized using Optuna. The optimization process specifically targeted maximizing the weighted F1 score. For strict reproducibility, the final parameter combinations for the \textit{Wine} and \textit{Credit} datasets are documented in \Cref{tab:rf_params}.

\begin{table}[htbp]
\centering
\caption{Optimized hyperparameters for the Random Forest classifier obtained via Optuna optimization.}
\label{tab:rf_params}
\begin{tabular}{lcc}
\toprule
Parameter & \textit{Wine} & \textit{Credit} \\
\midrule
\path{n_estimators} & 221 & 381 \\
\path{max_depth} & 5 & 9 \\
\path{min_samples_split} & 8 & 10 \\
\path{min_samples_leaf} & 3 & 2 \\
\path{max_features} & \path{sqrt} & \path{None} \\
\bottomrule
\end{tabular}
\end{table}

\todo[inline]{hier noch die korrekten werte nach dem finalen run eintragen}

\subsection{XGBoost}
\label{implementation_xgb}

The next tree-based ensemble model is the XGBoost classifier. This model serves as a direct point of comparison to the Random Forest. A primary reason for including XGBoost in this study is that it allows for a direct comparison of computational efficiency between CPU and GPU processing. To achieve this, the configuration dictionaries are kept identical for both executions. The only difference is the hardware device parameter, which is set to use the GPU. This controlled setup ensures a precise evaluation of the energy and emissions generated by the exact same model architecture on different hardware. This setup directly addresses the third research question: \textit{How does energy consumption differ between CPU-based and GPU-based training for equivalent tasks?}.

Additionally, XGBoost has a specific technical requirement for handling target variables. While the Random Forest algorithm natively adapts to different class structures, XGBoost requires explicit instructions for its learning objective. As mentioned in \Cref{experimental_dataset}, the \textit{Wine} dataset consists of three distinct categories. Therefore, the objective parameter is set to \texttt{multi:softmax}, paired with the \texttt{mlogloss} evaluation metric and a specified number of classes. In contrast, the \textit{Credit} and \textit{Higgs} datasets contain only two categories. Consequently, their objective is configured as \texttt{binary:logistic} and evaluated using the standard \texttt{logloss} metric.

Similar to the previous baseline algorithm, the parameters tailored to each dataset were determined using the Optuna optimization framework. Furthermore, the algorithm processes the unscaled feature matrix $X$ directly. It also uses the globally defined random state parameter to guarantee a deterministic execution across all hardware setups. The resulting configurations used for the final training phase are detailed in Table \ref{tab:xgb_params}.

\begin{table}[htbp]
\centering
\caption{Optimized hyperparameters for the XGBoost classifier across all datasets.}
\label{tab:xgb_params}
\begin{tabular}{lccc}
\toprule
Parameter & \textit{Wine} & \textit{Credit} & \textit{Higgs} \\
\midrule
\path{objective} & \path{multi:softmax} & \path{binary:logistic} & \path{binary:logistic} \\
\path{eval_metric} & \path{mlogloss} & \path{logloss} & \path{logloss} \\
\path{num_class} & 3 & N/A & N/A \\
\path{n_estimators} & 464 & 108 & 475 \\
\path{max_depth} & 7 & 3 & 8 \\
\path{learning_rate} & 0.0139 & 0.2878 & 0.2428 \\
\path{subsample} & 0.7790 & 0.9640 & 0.7100 \\
\path{colsample_bytree} & 0.7500 & 0.8430 & 0.9368 \\
\path{min_child_weight} & 5 & 7 & 2 \\
\bottomrule
\end{tabular}
\end{table}


\subsection{Multilayer Perceptron}
\label{implementation_mlp}

As discussed in \Cref{design_models}, a Multilayer Perceptron was selected to represent the neural network family. The implementation uses PyTorch as the primary deep-learning framework, combined with the Skorch library. Skorch wraps the PyTorch model in an interface that is fully compatible with the scikit-learn ecosystem. This is necessary because the entire evaluation logic, including cross-validation and pipeline construction, relies heavily on scikit-learn.

A dynamic architecture was chosen to ensure the network structure directly reflects the optimal parameters found during the Optuna tuning process. At runtime, the script reads the best parameter set from a JSON file generated during the tuning phase. It extracts the number of hidden layers via \texttt{n\_layers} and the neuron count for each individual layer via \texttt{layer\_i}. These values are then assembled into a list called \texttt{layer\_sizes}. This list is subsequently passed to the \texttt{MLPModule} class. This class constructs the full network dynamically upon instantiation, making the architecture entirely data-driven.

The \texttt{MLPModule} inherits from \texttt{nn.Module} and builds the hidden layers by iterating over the previously mentioned list. Each hidden block consists of four sequential operations. It starts with a fully connected linear transformation using \texttt{nn.Linear}, followed by a batch normalization step using \texttt{nn.BatchNorm1d}, a ReLU activation function, and a final dropout layer. Batch normalization normalizes the activations across the current mini-batch. This step stabilizes the training process and reduces internal covariate shifts \cite{ioffeBatchNormalizationAccelerating2015a}. The dropout layer randomly deactivates a fraction of neurons during the training phase. It serves as a regularization technique to prevent the network from memorizing the data \cite{srivastavaDropoutSimpleWay2014}. After processing all hidden blocks, a single linear output layer maps the final representations to raw class logits. The dimension of this output equals the exact number of target classes. No activation function is applied to this final layer, because the \texttt{CrossEntropyLoss} function explicitly expects unnormalized inputs.

\todo[inline]{das alles ist ziemlich theoretisch und sollte in die 02 eher rein}

To minimize the overall training time, the implementation checks for CUDA availability using \texttt{torch.cuda.is\_available()}. The algorithm runs on the GPU if one is present and automatically falls back to the CPU otherwise. PyTorch requires highly specific numeric types, and neural networks react sensitively to implicit type mismatches. Therefore, the feature matrix $X$ is explicitly cast to \texttt{float32} and the target vector $y$ to \texttt{int64} before the training begins.

The neural network is embedded within a scikit-learn \texttt{Pipeline} together with a \texttt{StandardScaler}. Since Multilayer Perceptrons lack built-in scale invariance, standardizing the inputs to a mean of zero and a unit variance is necessary to keep the gradient flow stable. The \texttt{NeuralNetClassifier} provided by Skorch receives all architecture parameters through a specific naming convention. It uses the \texttt{module\_\_*} prefix to forward keyword arguments directly to the constructor of the \texttt{MLPModule}. The Adam optimizer and the \texttt{CrossEntropyLoss} criterion are used consistently across all datasets. The learning rate and the batch size are taken directly from the tuning results and are therefore tailored to each dataset. A fixed random seed is applied using \texttt{torch.manual\_seed} to guarantee exact reproducibility across all iterations.

The maximum number of training epochs is set to 200. This value acts as an upper boundary rather than a strict target. An early stopping callback (\path{skorch.callbacks.EarlyStopping}) monitors the validation loss after each individual epoch. It halts the training process immediately if no improvement is observed for 10 consecutive epochs. This mechanism prevents unnecessarily long training runs and provides an additional layer of protection against overfitting. Consequently, the upper limit of 200 epochs is rarely reached in practice.

\todo[inline]{hier vielleicht eine visualisierung der layers einfügen}
\section{Hyperparameter Optimization}
\label{implementation_tuning}

\subsection{Optuna as Tuning Framework}
\label{implementation_optuna}

\todo[inline]{Nur Implementation, keine Begründung (die steht in \Cref{design_hyperparameter}): 

optuna.create\_study(direction="maximize"). 

TPESampler(seed=RANDOM\_STATE) für MLP (reproduzierbare Trial-Reihenfolge). study.optimize(objective, n\_trials=50). Tuning-Ergebnisse werden direkt als Hyperparameter-Dicts in die jeweiligen Model-Scripts übernommen.}

\subsection{Search Spaces per Model}
\label{implementation_search_spaces}

\todo[inline]{Suchräume für XGB (n\_estimators, max\_depth, learning\_rate, subsample, colsample\_bytree, min\_child\_weight), RF (n\_estimators, max\_depth, min\_samples\_split, min\_samples\_leaf, max\_features), MLP (n\_layers, layer\_sizes, dropout\_rate, lr, batch\_size). Tabelle oder kompakte Auflistung.}

\subsection{Tuning Strategy for HIGGS}
\label{implementation_tuning_higgs}

\todo[inline]{MLP-Tuning auf HIGGS: stratifiziertes Subset via --tune-sample-size statt vollem Dataset. Train-Test-Split (80/20) statt CV für Tuning (Begründung: Laufzeit). Ergebnisse in best\_params\_mlp\_higgs.json.}

\section{Emissions Measurement}
\label{implementation_emissions}

Energy consumption and CO\textsubscript{2} emissions were tracked using \textit{CodeCarbon} as described in \Cref{sec:theoretical_codecarbon}. For the hardware configuration used in this study, GPU power was measured directly via hardware counters, while RAM was estimated at 10W. CPU power could not be measured directly on Windows and was therefore estimated at 32.5W based on 50\% of the AMD Ryzen 7 5700X TDP, as discussed in \Cref{sec:theoretical_codecarbon}. \Cref{fig:emissions_flow} shows a visual representation of the process, from pre-processing to saving the results. Prior to training, data is loaded from the respective
dataset and, where required, features are normalized using
\texttt{StandardScaler}. Specifically, scaling was applied to Logistic
Regression and MLP, as these models are sensitive to the magnitude of input
features. Tree-based models (Random Forest, XGBoost) do not require
feature scaling and were trained on the raw input data.

Each training run was wrapped in an \texttt{EmissionsTracker} instance, which
monitored energy consumption across the full 5-fold cross-validation procedure.
Upon completion, the tracker was
stopped and the resulting emissions value, alongside accuracy, weighted F1-score,
training time, and dataset size, was saved to a central \texttt{results.csv}
file via the \texttt{save\_results()} utility function. In addition,
\textit{CodeCarbon} writes its own emissions log to the \texttt{emissions/}
directory, providing a detailed record of all tracked runs. The same tracking
approach was applied during hyperparameter tuning, ensuring that the energy
cost of the tuning process is also accounted for in the overall energy budget
of this study.

\begin{figure}[H]
\centering
\begin{tikzpicture}[
    node distance=1cm,
    box/.style={rectangle, rounded corners, draw=black,
                minimum width=6cm, minimum height=1cm,
                text centered, font=\small},
    tracker/.style={rectangle, rounded corners, draw=black,
                    dashed, minimum width=7cm, minimum height=0.8cm,
                    text centered, font=\small, fill=gray!10},
    arrow/.style={-{Stealth}, thick}
]

\node[box] (load)    {Load dataset};
\node[box, below=of load] (prep)    {Preprocess / scale features};
\node[tracker, below=of prep] (start)   {Start \texttt{EmissionsTracker}};
\node[box, below=of start] (cv)      {5-fold cross-validation};
\node[box, below=of cv] (train)   {Train \& evaluate model per fold};
\node[tracker, below=of train] (stop)    {Stop \texttt{EmissionsTracker}};
\node[box, below=of stop] (save)    {Save results to \texttt{results.csv}};

\draw[arrow] (load)  -- (prep);
\draw[arrow] (prep)  -- (start);
\draw[arrow] (start) -- (cv);
\draw[arrow] (cv)    -- (train);
\draw[arrow] (train) -- (stop);
\draw[arrow] (stop)  -- (save);

\end{tikzpicture}
\caption{Training and emissions tracking procedure per model-dataset combination.
         Dashed boxes indicate CodeCarbon \texttt{EmissionsTracker} calls.}
\label{fig:emissions_flow}
\end{figure}

\todo[inline]{Ggf. noch Subsections ergänzen: 5.6.1 CodeCarbon Integration, 5.6.2 Measurement Granularity (pro Modell × Dataset, eindeutige project\_names), 5.6.3 Limitations (CPU-Schätzung auf Windows, Background-Prozesse)}

\section{Inference Latency Measurement}
\label{implementation_inference}

Inference time was measured separately after training, outside the emissions
tracker. For each model, a single sample from the test set was passed through
the trained model and the elapsed time was recorded using
\texttt{time.perf\_counter()}. This provides a hardware-level estimate of
per-sample prediction latency, which is relevant in deployment scenarios where
repeated inference can accumulate significant energy costs
\cite{desislavovTrendsAIInference2023}.

\todo[inline]{Ergänzen: Begründung Single-Row statt Batch (Isolation von Training vs. Inference, kein Batch-Effekt). Modell kommt aus erstem CV-Fold (cv\_results["estimator"][0]).}

\section{Results Persistence and Orchestration}
\label{implementation_results}

\subsection{Results Schema}
\label{implementation_results_schema}

\todo[inline]{results.csv Schema beschreiben: timestamp, model, dataset, nrows, accuracy, f1, co2eq\_kg, training\_time\_s. Append-Logik (kein Überschreiben) — Hinweis auf mögliche Duplikate bei Mehrfachläufen. inference\_time.csv analog. csv.DictWriter, Header-Check.}

\subsection{Experimental Orchestration}
\label{implementation_orchestration}

\todo[inline]{run\_models.py: subprocess.run() pro Dataset × Modell-Kombination. Prozess-Isolation wichtig für CodeCarbon (prozessbasierte Messung). Fehlertoleranz: einzelne fehlgeschlagene Scripts blockieren nicht den Gesamtlauf.}
